\begin{explanationbox}
	\textbf{\textcolor{CExplanation}{一句话直觉}}
	函数图像整体上下平移时，纵坐标按平移量加减，横坐标不变。

	\medskip
	\textbf{\textcolor{CExplanation}{核心拆解}}
	\begin{description}[style=nextline, leftmargin=2.8em, labelsep=0.5em, itemsep=0.45em]
		\item[\textbf{要点一（平移方向）：}]
			向上平移后解析式变为$y=f(x)+b$，向下平移后变为$y=f(x)-b$。口诀“上加下减”。
		\item[\textbf{要点二（坐标变化）：}]
			平移只改变纵坐标，横坐标保持不变。变换前后对应点的横坐标相同。
	\end{description}

	\medskip
	\textbf{\textcolor{CExplanation}{几何本质（沿y轴平移）}}
	将函数图像上所有点同时沿$y$轴方向移动相同的距离$b$，图像整体在竖直方向上移动，形状不变。

	\medskip
	\textbf{\textcolor{CExplanation}{代数意义（函数值调整）}}
	变换前后，对同一个$x$，函数值$f(x)$变成$f(x)+b$或$f(x)-b$，即函数值整体增加或减少一个常数$b$。

	\medskip
	\textbf{\textcolor{CExplanation}{考点价值（正逆应用）}}
	\begin{itemize}[leftmargin=*, itemsep=0.5em, topsep=0.4em]
		\item \textbf{考法一（正向平移）：} 给出原函数和平移方向与距离，直接写出平移后的解析式。套用“上加下减”即可。
		\item \textbf{考法二（逆向求参）：} 已知平移后的解析式，推断平移量。例如，若$y=f(x)+k$是由$y=f(x)$平移得到，则$k>0$表示上移$|k|$，$k<0$表示下移$|k|$。
	\end{itemize}

	\medskip
	\textbf{\textcolor{CExplanation}{顿悟点}}
	牢记“上加下减”专指$y$轴方向平移，与$x$轴方向的“左加右减”区分。加减是直接作用在函数值$f(x)$上，而非$x$本身。

	\medskip
	\textbf{\textcolor{CExplanation}{使用场景}}
	\begin{itemize}[leftmargin=*, itemsep=0.5em, topsep=0.4em]
		\item \textbf{场景一（图像变换作图）：} 已知$y=f(x)$的图像，要得到$y=f(x)+b$的图像，只需将原图像向上平移$b$个单位。
		\item \textbf{场景二（解析式识别）：} 见到$y=2x^2+1$，可看作$y=2x^2$向上平移1个单位。
	\end{itemize}

\end{explanationbox}
